% =============================================================
% 第5章 技術性能指標TPMの推移 (課題8) 章本体
% reports.tex から \input される．
% =============================================================

\section{はじめに}
本章には，Arduinoオーケストラ開発における技術性能指標（TPM: Technical Performance Measure）の測定推移を記録する．
TPM項目の定義はマネジメント計画書（\texttt{management\_plan.tex}）の表「TPM一覧」を参照．

\section{TPM項目（参考）}

\begin{table}[H]
  \centering
  \begin{tabularx}{\linewidth}{lX}
    \toprule
    No. & TPM（監視値・許容範囲） \\
    \midrule
    TPM-1 & 距離計算誤差：$\leq \pm 3\,\mathrm{cm}$ \\
    TPM-2 & 拍検出成功率：$\geq 95\,\%$ \\
    TPM-3 & Arduino→Processing転送遅延：$\leq 50\,\mathrm{ms}$ \\
    TPM-4 & EMA（$\alpha=0.2$）テンポ収束拍数：$\leq 13$拍 \\
    TPM-5 & 指揮者人形が追従できる最大BPM：$\geq 104\,\mathrm{bpm}$ \\
    TPM-6 & テンポ予測誤差（予測拍時刻と実測拍時刻の差）：$\leq 1\,\mathrm{ms}$ \\
    \bottomrule
  \end{tabularx}
\end{table}

\section{TPM測定推移}

\begin{longtable}{p{2cm}p{1.5cm}p{2.5cm}p{2.5cm}p{1.5cm}p{4cm}}
\toprule
測定日 & TPM & 目標値 & 実測値 & 達否 & 備考 \\
\midrule
\endfirsthead

\toprule
測定日 & TPM & 目標値 & 実測値 & 達否 & 備考 \\
6月10日 & & & & & まだ性能テストする段階まで達しておらず，TPMは達成できていない．\\
\midrule
\endhead

% --- 各TPMの測定推移を記載 ---
% 書式: YYYY-MM-DD & TPM-X & 目標値 & 実測値 & ○/× & 備考 \\
2026-06-17 & TPM-6 & $\leq 1\,\mathrm{ms}$ & $0.025\,\mathrm{ms}$ & ○ & 100回検出による平均誤差（UNIT-4単体テスト時に計測） \\

\bottomrule
\endfoot
\end{longtable}

\section{TPM推移グラフ（任意）}

% TPMごとの推移を折れ線グラフで可視化する場合は tpm_history/figures/ に画像を配置．
%
% \begin{figure}[H]
%   \centering
%   \includegraphics[width=0.8\linewidth]{tpm_history/figures/tpm_chart.pdf}
%   \caption{TPM推移グラフ}
% \end{figure}
